\documentclass[11pt]{article}

% ============================================================
%  S, T, U in warped (Randall--Sundrum) models
%  A slow, pedagogical internal review.
%  Written to track (a) the physics, (b) what is in OUR code,
%  (c) the subtleties and gaps.
%  This is an internal working note, not for distribution.
% ============================================================

\usepackage[margin=1in]{geometry}
\usepackage{amsmath,amssymb,amsthm}
\usepackage{mathtools}
\usepackage{graphicx}
\usepackage{booktabs}
\usepackage{enumitem}
\usepackage{xcolor}
\usepackage{framed}
\usepackage[colorlinks=true,linkcolor=blue,citecolor=blue,urlcolor=blue]{hyperref}

% lightweight "callout" box
\colorlet{shadecolor}{gray!12}
\newenvironment{callout}[1]{%
  \par\medskip\begin{shaded}\noindent\textbf{#1}\par\smallskip}{\end{shaded}\medskip}

\newcommand{\MKK}{M_{\mathrm{KK}}}
\newcommand{\LIR}{\Lambda_{\mathrm{IR}}}
\newcommand{\sw}{s_W}
\newcommand{\cw}{c_W}
\newcommand{\swsq}{s_W^2}
\newcommand{\cwsq}{c_W^2}
\newcommand{\Pibig}{\Pi}
\newcommand{\half}{\tfrac{1}{2}}
\newcommand{\vev}{v}
\newcommand{\eps}{\varepsilon}
\newcommand{\SUL}{SU(2)_L}
\newcommand{\SUR}{SU(2)_R}
\newcommand{\UX}{U(1)_X}
\newcommand{\PLR}{P_{LR}}
\newcommand{\Tparam}{\hat{T}}
\newcommand{\code}[1]{\texttt{#1}}
\newcommand{\todo}[1]{{\color{red}[#1]}}

\theoremstyle{definition}
\newtheorem{definition}{Definition}

\title{\bf Oblique Electroweak Parameters $S$, $T$, $U$ in Warped Models\\[4pt]
\large A slow, pedagogical internal review --- physics, code, and subtleties}
\author{Internal working note (5D--Neutrino--Mixing repo)}
\date{\today}

\begin{document}
\maketitle

\begin{abstract}
\noindent
This note is a from-scratch, deliberately slow walk-through of the
Peskin--Takeuchi oblique parameters $S$, $T$, $U$ and how they constrain a
Randall--Sundrum (RS) warped extra dimension --- the single least-rigorous
ingredient in our flavour/EW constraint catalogue. It assumes no prior
knowledge of ``oblique parameters'' or ``custodial symmetry'': both are built
up from the beginning. Part~\ref{part:sm} defines the framework in the Standard
Model (SM). Part~\ref{part:rs} explains \emph{why} warped models generate a
dangerous $T$ and a positive $S$, and where the famous volume-log enhancement
comes from. Part~\ref{part:cust} explains the custodial cure
($\SUL\times\SUR\times\UX\times\PLR$), the separate $\PLR$ protection of
$Z\bar b_L b_L$, and the calculable one-loop top-partner corrections.
Part~\ref{part:code} documents exactly what our code does
(\code{quarkConstraints/oblique\_stu.py}, \code{EW001}), and
Part~\ref{part:subtle} catalogues the subtleties and the precise gap between
``what we have'' (a leading-order proxy) and ``rigorous.'' An annotated reading
guide closes the note.
\end{abstract}

\tableofcontents

% ====================================================================
\part*{Orientation: why this note exists}
\addcontentsline{toc}{section}{Orientation: why this note exists}
% ====================================================================

In our constraint catalogue the flavour-changing observables
($\eps_K$, $\Delta m_d$, $\Delta m_s$, $D^0$--$\bar D^0$, $\mu\to e\gamma$) and
$Z\to b\bar b$ are computed at a level we are willing to call
\emph{rigorous}. The oblique electroweak parameters are \emph{not}: \code{EW001}
uses a closed-form parametric \emph{proxy}, not a per-point spectral
computation. Yet $S$ and $T$ are physically among the most important
constraints on any warped model: in the minimal (non-custodial) version the
$T$ parameter alone pushes the Kaluza--Klein (KK) scale to
$\sim 8$--$10$~TeV, and after one adopts custodial protection it is the $S$
parameter that becomes the binding electroweak constraint. To know what our
scan is really doing --- and what it would take to make it rigorous --- we need
to understand $S$, $T$, $U$ from the ground up. That is the purpose of this
note.

A one-paragraph executive summary, to be unpacked over the following pages:

\begin{quote}
\small
New heavy physics that couples to the SM mainly through the electroweak gauge
bosons shows up, to leading order, as just a few numbers --- the
\emph{oblique parameters} $S,T,U$ --- which measure shifts in the $W$ and $Z$
vacuum polarisations. $T$ measures the breaking of \emph{custodial
isospin} (the symmetry that protects $M_W = M_Z\cw$ at tree level); $S$
measures the overall size of the new electroweak sector. In RS, the Higgs lives
on (or near) the IR brane and the KK gauge bosons are localised there too, so
integrating them out generates $S\sim \vev^2/\MKK^2$ and, crucially, a $T$ that
is \emph{enhanced by the large volume factor} $L=k\pi r_c \approx 35$ because the
minimal bulk gauge group has no custodial symmetry. Enlarging the bulk gauge
group to $\SUL\times\SUR\times\UX$ supplies a custodial $\SUR$ that removes the
volume enhancement of $T$ (flipping its sign and suppressing it by $1/L^2$),
while a discrete $L\!\leftrightarrow\! R$ parity $\PLR$ separately protects the
$Z\bar b_L b_L$ vertex. $S$ is \emph{not} protected and sets the residual
$\sim$few-TeV electroweak floor.
\end{quote}

\noindent The logical spine of the whole note, in one picture:
\begin{footnotesize}
\begin{verbatim}
  custodial isospin  SU(2)_L x SU(2)_R --> SU(2)_V   ==guards==>  rho = 1  ==>  T = 0
                           |
  minimal RS bulk = SU(2)_L x U(1)_Y   == NO custodial ==>
                           |              Delta T ~ (pi L / 2cw^2)(v^2/M^2),  L = ln(M_Pl/TeV) ~ 35
                           |              ==> volume-enhanced T  ==> KK floor ~ 8-10 TeV
                           |
  cure: enlarge bulk gauge group to  SU(2)_L x SU(2)_R x U(1)_X
       |
       +-- SU(2)_R  (continuous)  ==> kills volume-enhanced T   (coeff: +pi L/2 --> -pi/4L)
       |
       +-- P_LR  (discrete L<->R, needs bidoublet)  ==> protects Z b_L b_L   [SEPARATE from T!]
                           |
  S parameter: protected by NEITHER  ==>  residual EW floor ~ 3 TeV  (binding once T is cured)
\end{verbatim}
\end{footnotesize}
\noindent Keep this map in mind; every box in it is unpacked below. The most
common confusion --- that ``custodial'' is one knob that fixes everything --- is
already defused here: there are \emph{two} distinct arrows ($\SUR$ for $T$,
$\PLR$ for $Zb_L$), and $S$ is protected by neither.

% ====================================================================
\part{The Standard-Model framework: what $S$, $T$, $U$ mean}
\label{part:sm}
% ====================================================================

\section{Oblique vs.\ non-oblique: the key assumption}
\label{sec:oblique}

Imagine some new physics at a scale $M \gg M_Z$. Generically it corrects every
electroweak observable: $Z$-pole asymmetries, the $W$ mass, $\Gamma_Z$, atomic
parity violation, and so on. Cataloguing all of these correction-by-correction
would need many parameters. There is, however, a very common and very useful
special case.

\begin{definition}[Oblique corrections]
New physics is called \emph{oblique} (or \emph{universal}) if it enters only
through the \emph{self-energies} (vacuum polarisations) of the electroweak gauge
bosons $\gamma, W^\pm, Z$ --- i.e.\ through corrections to the gauge-boson
propagators --- and \emph{not} through direct vertex corrections to specific
fermions or through box diagrams. Equivalently: the new physics couples to SM
fermions only via the conserved electroweak currents.
\end{definition}

Why is this so often a good approximation? Two reasons. First, if the new
states are heavy colour/electroweak multiplets that mix with gauge bosons but
couple to light fermions only through gauge interactions, their leading effect
\emph{is} a propagator correction. Second, the dominant SM radiative corrections
(from the top quark and the Higgs) are themselves oblique, so the formalism was
historically built to package them.

The power of the oblique assumption is that the gauge self-energies, expanded to
linear order in $q^2$, are captured by a \emph{finite number of constants}. Let
$\Pibig_{XY}(q^2)$ denote the transverse part of the $\langle J_X J_Y\rangle$
current--current correlator, with $X,Y \in \{1,2,3,Q\}$ labelling the $\SUL$
generators $1,2,3$ and the electromagnetic current $Q$. Expand
\begin{equation}
\Pibig_{XY}(q^2) = \Pibig_{XY}(0) + q^2\, \Pibig'_{XY}(0) + \mathcal{O}(q^4).
\label{eq:Piexpand}
\end{equation}
Electromagnetic gauge invariance forces $\Pibig_{QQ}(0)=\Pibig_{3Q}(0)=0$ (the
photon stays massless), so the non-trivial constants at this order are
$\Pibig_{11}(0), \Pibig_{33}(0)$ and the derivatives
$\Pibig'_{11},\Pibig'_{33},\Pibig'_{3Q},\Pibig'_{QQ}$. After using three of
these to renormalise the three SM input parameters
($\alpha$, $G_F$, $M_Z$), exactly \emph{three} independent combinations remain.
Those are $S$, $T$, $U$.

\paragraph{Caveat we will return to.} The expansion \eqref{eq:Piexpand} truncates
at $\mathcal{O}(q^2)$. For new physics whose mass is not enormously above $M_Z$
--- precisely the light-KK regime we care about --- the $\mathcal{O}(q^4)$ terms
(the parameters usually called $W$ and $Y$) are not automatically negligible,
and RS in particular has genuinely \emph{non-oblique} pieces from
generation-dependent fermion localisation. We flag this here and develop it in
\S\ref{sec:nonoblique}.

\section{Custodial symmetry and the $\rho$ parameter (the heart of $T$)}
\label{sec:custodialSM}

Before defining $T$ formally it pays to understand the symmetry it measures,
because that same symmetry is the whole story in RS.

The SM Higgs is a complex doublet $H$, which has four real components. Write them
as a $2\times 2$ matrix
\begin{equation}
\Sigma = \big(\,\tilde H\ \ H\,\big), \qquad \tilde H = i\sigma_2 H^*,
\end{equation}
on which a \emph{global} $\SUL\times\SUR$ acts as
$\Sigma \to U_L\,\Sigma\,U_R^\dagger$. The Higgs potential depends only on
$\mathrm{tr}(\Sigma^\dagger\Sigma)=2|H|^2$, so it is invariant under this full
$\SUL\times\SUR$. When $H$ gets a vev, $\langle\Sigma\rangle = \frac{\vev}{\sqrt2}\,
\mathbb{1}$, the symmetry breaks to the diagonal subgroup
$\SUL\times\SUR \to SU(2)_V$. This unbroken \emph{vectorial} $SU(2)_V$ is called
\textbf{custodial isospin}.

Its job is to ``guard'' a relation. The three would-be Goldstone bosons (eaten by
$W^\pm, Z$) form a triplet of $SU(2)_V$, so their mass matrix is forced to be
proportional to the identity in isospin space. In the limit of exact custodial
symmetry (i.e.\ ignoring hypercoupling $g'$ and Yukawa splittings), this means
\begin{equation}
\rho \equiv \frac{M_W^2}{M_Z^2\,\cwsq} = 1
\label{eq:rho}
\end{equation}
at tree level, where $\cwsq = 1-\swsq$ and $\sw$ is the weak mixing angle. The
custodial symmetry is only \emph{approximate} in the SM --- it is broken
explicitly by (i) gauging hypercharge ($g'\neq 0$: in the SM $\SUR$ is only a
\emph{global} symmetry, and switching on $g'$ singles out its $T^3_R$ direction,
breaking $SU(2)_V$ explicitly) and (ii) the large top--bottom Yukawa splitting.
(In the warped \emph{custodial} construction of Part~\ref{part:cust} this becomes
sharp: $\SUR$ is gauged and hypercharge is $Y = T^3_R + (B\!-\!L)/2$.) These are exactly the
sources that make the SM $\rho$ deviate from 1 at loop level. The deviation is
the $T$ parameter:
\begin{equation}
\alpha\,T \;=\; \rho - 1 \;=\; \Delta\rho .
\label{eq:Tisrho}
\end{equation}
\textbf{Slogan to remember:} \emph{$T$ is a direct measure of custodial-isospin
breaking in the new sector.} A model that respects custodial symmetry
automatically has small $T$; a model that breaks it badly has a large $T$. RS is
the second kind unless one engineers the first.

\section{Formal definitions of $S$, $T$, $U$}
\label{sec:STUdef}

With the physical meaning in hand, the Peskin--Takeuchi definitions are:
\begin{align}
\alpha\,T &= \frac{\Pibig_{11}(0) - \Pibig_{33}(0)}{M_W^2},
\label{eq:Tdef}\\[4pt]
\alpha\,S &= 4\swsq\cwsq \,\big[\Pibig'_{33}(0) - \Pibig'_{3Q}(0)\big],
\label{eq:Sdef}\\[4pt]
\alpha\,U &= 4\swsq\,\big[\Pibig'_{11}(0) - \Pibig'_{33}(0)\big].
\label{eq:Udef}
\end{align}
Read these physically:
\begin{itemize}[leftmargin=1.4em]
\item \textbf{$T$} \eqref{eq:Tdef} is the \emph{difference} between the charged
($\Pibig_{11}$) and neutral ($\Pibig_{33}$) current self-energies at zero
momentum --- an isospin-breaking ``mass-splitting'' of the $W$ and $Z$ sectors.
Exactly custodial $\Rightarrow$ $\Pibig_{11}(0)=\Pibig_{33}(0)$ $\Rightarrow$ $T=0$.
It is dimensionless because of the $1/M_W^2$.
\item \textbf{$S$} \eqref{eq:Sdef} measures the \emph{momentum-dependence}
(the $\Pibig'$) of the neutral-current self-energy --- intuitively, the
``size'' or number of new electroweak-charged states, independent of whether
they break custodial symmetry. A heavy fourth generation, a technicolor sector,
or a tower of KK modes all contribute positively.
\item \textbf{$U$} \eqref{eq:Udef} is a second isospin-breaking combination but
involves \emph{derivatives}; it is generically of order $(M_Z/M_{\rm NP})^2$
smaller than $S$ and $T$,\footnote{In the EFT-operator language $S$ and $T$ are
generated by dimension-6 operators ($\mathcal{O}_{WB}$ and
$|H^\dagger D_\mu H|^2$), whereas $U$ first arises at dimension 8 --- the same
fact, stated as operator counting. The suppression argument in the text is the
physical reason.} so it is almost always negligible and is routinely set to
$U=0$ in fits. Our code does this.
\end{itemize}

By construction all three vanish in the SM with a reference Higgs and top mass:
$S=T=U=0$ is the SM point, and $S,T,U$ measure \emph{deviations} from it.

\section{The experimental constraint: a correlated ellipse}
\label{sec:expt}

Electroweak data constrain $(S,T)$ (with $U=0$) to a tilted error ellipse. The
values our catalogue uses (PDG~2025 global fit, $U$ fixed), stored in
\code{EW001.yaml}, are
\begin{equation}
S = 0.026 \pm 0.075, \qquad T = 0.047 \pm 0.066, \qquad \rho_{ST} = +0.90,
\label{eq:fit}
\end{equation}
with a strong positive correlation $\rho_{ST}=0.90$. ``$U$ fixed'' means $U\equiv0$
is imposed in the fit, which tightens the $S$--$T$ errors relative to the
three-parameter floating fit (the catalogue also stores the floating values
$S=0.021\pm0.096$, $T=0.04\pm0.12$, $U=0.008\pm0.092$ for reference). A model
point with prediction $(S_p,T_p)$ is tested with the Gaussian $\chi^2$
\begin{equation}
\chi^2 = (\Delta x)^\top \,C^{-1}\, (\Delta x), \qquad
\Delta x = \begin{pmatrix} S_p - 0.026 \\ T_p - 0.047 \end{pmatrix},\quad
C = \begin{pmatrix} \sigma_S^2 & \rho_{ST}\sigma_S\sigma_T \\
\rho_{ST}\sigma_S\sigma_T & \sigma_T^2 \end{pmatrix},
\label{eq:chi2}
\end{equation}
and ``passes'' at 95\% C.L.\ for two degrees of freedom if
$\chi^2 \le 5.991$. This is exactly what \code{compare\_oblique\_st\_to\_fit}
implements (\S\ref{part:code}); note the \emph{experimental} side of \code{EW001}
is genuinely rigorous --- it is only the \emph{prediction} $(S_p,T_p)$ that is a
proxy.

The strong positive correlation matters: because $\rho_{ST}=+0.90$, the
allowed region is a thin diagonal band. A model that raises $S$ and $T$
\emph{together} (along the band) is far less constrained than one that raises
only one of them. RS, unfortunately, tends to push $T$ hard while moving $S$
only modestly --- across the band, not along it:
\begin{footnotesize}
\begin{verbatim}
        T
        ^                         . . allowed 95% band (thin, tilted by rho=+0.90)
        |                    . .'        '. .
   +0.4 |                .'      o  <-- minimal RS:  pushes T up, ACROSS the band  (excluded)
        |            . .'      '. .
        |        . .'  *  <-- SM (0,0), data centre near here
        |    . .'      '. .
      0 +--.'------------'.------------------> S
        |      ^ along the band = raise S and T together (mildly constrained)
       -+------+---------------------------------
              0
\end{verbatim}
\end{footnotesize}
The lesson the picture encodes: a model that moves \emph{along} the diagonal band
is barely constrained; one that moves \emph{across} it (large $\Delta T$ at small
$\Delta S$, as minimal RS does) is excluded quickly.

% ====================================================================
\part{Why warped models generate $S$ and $T$}
\label{part:rs}
% ====================================================================

\section{RS recap, in just enough detail}
\label{sec:rsrecap}

Recall the geometry (consistent with the repo conventions). The metric is a slice
of AdS$_5$,
\begin{equation}
ds^2 = e^{-2k|y|}\eta_{\mu\nu}dx^\mu dx^\nu - dy^2, \qquad 0 \le y \le \pi r_c,
\end{equation}
with a UV brane at $y=0$ (Planck scale) and an IR brane at $y=\pi r_c$
(TeV scale). The warp factor is $\eps = e^{-k\pi r_c} \sim 10^{-15}$ and the all
important \textbf{volume factor} is
\begin{equation}
\boxed{\;L \equiv k\pi r_c = \ln(1/\eps) = \ln(M_{\rm Pl}/\mathrm{TeV}) \approx 35.\;}
\label{eq:L}
\end{equation}
This single large number is what makes the hierarchy problem soluble --- and, as
we are about to see, it is the same number that enhances the $T$ parameter. Our
code stores it as \code{DEFAULT\_RS\_VOLUME\_LOG = 35.0}.

A bulk gauge field has a tower of KK modes. The zero mode is the SM gauge boson,
flat in $y$ (up to the warp); the excited modes are localised toward the IR
brane and have masses
\begin{equation}
m_n \simeq \left(n - \tfrac14\right)\pi\,k\,\eps \equiv x_n\,\LIR ,
\label{eq:KKmass}
\end{equation}
where the $x_n$ are roots of a Bessel function. The right-hand side is the
\emph{large-$n$ asymptote}; for $n=1$ it gives $0.75\pi\approx 2.36$. The exact
first root depends on the boundary conditions: the $\eps\!\to\!0$ Neumann root is
the first zero of $J_0$, $x_1 = 2.405$, while the physical gauge mode with the
finite-$\eps$ (NN) boundary conditions used in our solver sits slightly higher,
$x_1 \approx 2.45$ (this is the value pinned in the repo's
\code{rs\_ew\_spectrum} derivation). The first KK gauge mass is thus
$\MKK^{\rm phys} = x_1 \LIR \approx 2.45\,\LIR$.
\textbf{This factor of $2.45$ is a recurring source of confusion} (see
\S\ref{sec:Mconvention}): ``$\MKK$'' can mean the physical first KK mass or the
scale $\LIR$ itself, and the two differ by $\sim 2.45$, which is a large factor
when quoting TeV bounds.

\section{The physical origin of the oblique corrections}
\label{sec:rsorigin}

Here is the mechanism in words. The Higgs is localised on (or very near) the IR
brane. The SM $W$ and $Z$ are the \emph{flat} zero modes of the bulk gauge
fields. When the Higgs gets its vev, it does two things:
\begin{enumerate}[leftmargin=1.6em]
\item it gives the zero modes their usual masses $M_W, M_Z$; and
\item because the vev sits at the IR brane where the KK modes pile up, it
\emph{mixes} the zero modes with the heavy KK modes.
\end{enumerate}
Integrating out the heavy KK modes then shifts the effective $W/Z$ propagators
--- precisely an oblique correction. Diagrammatically, the leading effect is a
mass-insertion mixing $\propto \vev^2$ followed by a heavy KK propagator
$\propto 1/\MKK^2$, so \emph{every} oblique parameter scales as
\begin{equation}
\{S, T\}\ \sim\ \frac{\vev^2}{\MKK^2}\,\times\,(\text{coefficient}),
\label{eq:scaling}
\end{equation}
which is the universal $\vev^2/\MKK^2$ decoupling. The whole game is in the
coefficient --- and the coefficients of $S$ and $T$ behave very differently.

There is an equivalent and very useful \emph{holographic / dual} picture
(Agashe--Delgado--May--Sundrum). A bulk gauge field is dual to a 4D CFT operator;
the IR-localised Higgs is dual to a strongly-coupled composite Higgs sector. In
that language, the minimal model's $T$ problem is simply the statement that the
composite Higgs sector \emph{lacks custodial isospin symmetry}. This is why the
fix is symmetry-based, not fine-tuning-based.

\section{The $T$ parameter problem: where the volume factor comes from}
\label{sec:Tproblem}

Consider the $T$ parameter, \eqref{eq:Tdef}. It is the difference of the charged
and neutral self-energies at zero momentum. In the minimal bulk gauge group
($\SUL\times U(1)_Y$ only) there is \emph{nothing} that forces these to be equal:
the $W$ and the $Z$ mix with their respective KK towers with different strengths,
governed by the hypercharge coupling, and the difference does not cancel.

The crucial point is the \emph{size} of the coefficient, and here is the cleanest
way to see where the volume factor $L$ comes from. A gauge zero mode is flat
\emph{as a 5D wavefunction}, but ``flat'' is deceptive: it must be
\emph{canonically normalised} by integrating $|\psi_0|^2$ over the warped extra
dimension, and the AdS warping stretches that proper measure by the volume
factor --- the integral runs over an effective length $\propto L = \ln(1/\eps)$.
So the normalised zero mode is suppressed by $1/\sqrt{L}$ over most of the bulk
but is \emph{relatively enhanced at the IR brane}, where the warp factor is
smallest. The Higgs lives exactly there, on the IR brane, so it samples the
zero mode precisely where the volume normalisation has piled up the
wavefunction. Working through the overlap integral, the coupling of the would-be
SM gauge boson to the IR-Higgs --- and hence the $W$/$Z$ mixing with the KK tower
that feeds $\Pibig_{11}(0)-\Pibig_{33}(0)$ --- carries an explicit factor of $L$.
\textbf{In one line:} \emph{the gauge boson is normalised over a volume $\sim L$,
but the Higgs samples it only at the IR brane, and that mismatch is the factor
$L$ in $\Delta T$.} It is the very same $L = \ln(M_{\rm Pl}/\mathrm{TeV})$ that
generates the fermion mass hierarchies and solves the gauge hierarchy
problem --- which is why the $T$ problem is so stubborn: removing it naively
would mean giving up the warping that motivates the whole construction. The
parametric result, which our code encodes directly, is
\begin{equation}
\boxed{\;\Delta T_{\rm minimal} \;\simeq\; \frac{\pi L}{2\,\cwsq}\,
\frac{\vev^2}{\MKK^2}.\;}
\label{eq:Tmin}
\end{equation}
The dangerous feature is the explicit factor $L\approx 35$. Numerically, with
$\vev = 246$~GeV and $\cwsq \approx 0.769$, the coefficient is
$\pi L/(2\cwsq) \approx 71.5$. To bring $\Delta T$ below its $\sim 0.1$
experimental ceiling one needs $(\vev/\MKK)^2 \lesssim 1.4\times10^{-3}$, i.e.\
$\MKK \gtrsim 6$--$10$~TeV depending on how the $S$ correlation is folded in. This
is the celebrated \textbf{RS $T$-parameter problem}: custodial breaking by the
minimal bulk gauge group, amplified by the hierarchy volume $L$, forces the KK
scale uncomfortably high.

Two sanity remarks. (i) The enhancement is \emph{logarithmic in the hierarchy} but
\emph{order-35 in practice} --- ``$L$ is large'' is not a small-parameter
expansion gone wrong, it is the same big log that solves the hierarchy. (ii) The
sign is positive: minimal RS predicts $T>0$, which the data mildly \emph{prefer}
(the fit centre is $T=+0.047$), so the tension is one of magnitude, not sign.

\section{The $S$ parameter: positive, order one, and not volume-protected}
\label{sec:Sparam}

$S$ \eqref{eq:Sdef} measures the momentum-dependence of the neutral self-energy
--- effectively the number and coupling of new EW-charged states, here the KK
gauge tower coupling to light (UV-localised) fermions. Summing the tree-level KK
exchange gives a positive
\begin{equation}
\Delta S \;\simeq\; c_S\,\frac{\vev^2}{\MKK^2}, \qquad c_S = \mathcal{O}(\text{few}\!-\!\text{tens}),
\label{eq:Smin}
\end{equation}
with no factor of $L$ multiplying it (and in some treatments a mild $1/L$
\emph{suppression} of the leading piece). The key structural facts:
\begin{itemize}[leftmargin=1.4em]
\item $\Delta S > 0$ and does \emph{not} carry the volume enhancement that
$\Delta T$ does, so in the minimal model $T$ is the dominant constraint.
\item $S$ is \emph{not} protected by custodial symmetry --- $SU(2)_V$ says
nothing about $\Pibig'_{33}-\Pibig'_{3Q}$. Hence \textbf{after} one cures $T$
with custodial symmetry, $S$ becomes the binding electroweak constraint.
\item The coefficient $c_S$ is genuinely model-dependent (it depends on the gauge
representation content and on UV-brane kinetic terms). Our code takes
$c_S = 30$ from the catalogue anchor (\S\ref{part:code}); this is a
deliberately conservative O(tens) choice and is one of the things a rigorous
treatment would replace with a spectral computation (\S\ref{part:subtle}).
\end{itemize}

\section{The non-oblique caveat: RS is not purely $S,T,U$}
\label{sec:nonoblique}

Everything above assumed the new physics is purely oblique. RS violates this in
two ways that an honest treatment must keep in mind.

\paragraph{(a) Generation-dependent vertices.} Light fermions are UV-localised and
heavy fermions IR-localised, so the KK gauge bosons couple \emph{non-universally}
to different generations. The third generation in particular (the IR-localised
$t_L, b_L$) gets large direct vertex corrections --- this is exactly the
$Z\bar b_L b_L$ problem, which is \emph{not} an oblique effect and is treated
separately in our code (constraints \code{T010}/\code{T011}). One cannot fold the
$Zb\bar b$ shift into $S,T,U$; it is a genuinely different operator.

\paragraph{(b) Higher-derivative oblique parameters $W$, $Y$.} The
$S,T,U$ formalism keeps only the \emph{first} derivative of the self-energies at
$q^2=0$. That truncation is justified when the new physics is much heavier than
$M_Z$ --- but light KK modes sit only a factor of a few above $M_Z$, so the
\emph{next} derivative (the $\mathcal{O}(q^4)$ terms dropped in
\eqref{eq:Piexpand}) is not negligible. Barbieri--Pomarol--Rattazzi--Strumia
package these into two further universal parameters $W$ and $Y$. In their
hatted normalisation (which strips the awkward $\alpha/\swsq$ factors), all four
parameters are the same parametric size,
\begin{equation}
\hat S,\ \hat T,\ W,\ Y \ \sim\ \frac{v^2}{\MKK^2},
\end{equation}
so $W,Y$ are \emph{not} parametrically smaller than $\hat S,\hat T$ for light KK
states. The slogan: \emph{for genuinely heavy new physics three numbers suffice;
for light KK modes you need (at least) four, and $Z$-pole data alone cannot
separate them --- you also need the high-energy LEP2 / off-pole data that
constrain $W,Y$.} A fully rigorous warped EW fit therefore uses
$\{\hat S,\hat T,W,Y\}$, or better fits directly to the $Z$-pole and $M_W$
observables, rather than projecting onto two numbers. \textbf{Our proxy keeps
only $S$ and $T$ with $U=0$ and drops $W,Y$ entirely}; this is a known and
documented limitation, not an oversight.

% ====================================================================
\part{The custodial cure}
\label{part:cust}
% ====================================================================

\section{The idea: a bulk custodial symmetry}
\label{sec:custidea}

The cure (Agashe--Delgado--May--Sundrum) is to make the bulk \emph{respect}
custodial isospin, so that the dangerous $\Delta T$ is forbidden by symmetry
rather than tuned away. Concretely, enlarge the bulk gauge group from the SM to
\begin{equation}
\boxed{\;G_{\rm bulk} = \SUL \times \SUR \times \UX,\qquad X = (B-L)/2,\;}
\label{eq:custgroup}
\end{equation}
and break it down to the SM on the UV brane (by boundary conditions) while
preserving a custodial $SU(2)_V$ in the IR/Higgs sector. The Higgs is now a
bidoublet $(2,2)$ of $\SUL\times\SUR$, so its vev leaves an unbroken diagonal
$SU(2)_V$ --- exactly the custodial symmetry of \S\ref{sec:custodialSM}, now
built into the 5D theory.

In the dual language: we have given the composite Higgs sector a global custodial
$\SUR$, so the strongly-coupled dynamics no longer breaks isospin at leading
order.

\section{How $\SUR$ removes the volume enhancement of $T$}
\label{sec:Tcure}

The sharpest way to state the mechanism --- the one Casagrande et al.\ make
central --- is a parameter count at the IR brane. In the minimal model the IR
boundary conditions fix what become the $W$ and $Z$ mass terms
\emph{independently}: there are effectively \emph{two} mass parameters, and
nothing relates them, so $\Pibig_{11}(0)\neq\Pibig_{33}(0)$ and $T$ is generated
(and then volume-enhanced). With custodial $\SUR$, the Higgs is a bidoublet and
its vev enters the IR boundary condition through a \emph{single} mass parameter
that acts symmetrically on the charged and neutral sectors. One parameter cannot
split the charged and neutral self-energies, so $\Pibig_{11}(0)-\Pibig_{33}(0)$
is forced to cancel at the volume-enhanced order:
\begin{center}
\begin{tabular}{lll}
\toprule
 & IR mass parameters & consequence \\
\midrule
minimal $\SUL\times U(1)_Y$ & two (independent $M_W$, $M_Z$ terms) & $\Pibig_{11}(0)\neq\Pibig_{33}(0)\Rightarrow T\propto L$ \\
custodial $\SUL\times\SUR\times\UX$ & one (single bidoublet vev) & $\Pibig_{11}(0)=\Pibig_{33}(0)$ at leading order $\Rightarrow T$ tiny \\
\bottomrule
\end{tabular}
\end{center}
What survives is parametrically smaller. The structural
replacement that our code encodes is the coefficient flip
\begin{equation}
\boxed{\;c_T:\quad \frac{\pi L}{2\,\cwsq}\ \xrightarrow{\ \text{custodial}\ }\
-\,\frac{\pi}{4\,\cwsq\,L}.\;}
\label{eq:Tcust}
\end{equation}
Note three things about \eqref{eq:Tcust}:
\begin{enumerate}[leftmargin=1.6em]
\item \textbf{The $L$ moves from numerator to denominator}: the
$\mathcal{O}(L)$ \emph{enhancement} becomes an $\mathcal{O}(1/L)$
\emph{suppression}. Net change in the coefficient is a factor
$\sim 1/(2L^2)\sim 4\times10^{-4}$ --- $T$ is rendered essentially harmless.
\item \textbf{What survives is a small, model-dependent residual.} The leading
custodial-symmetric piece cancels; the leftover is parametrically
$\mathcal{O}(1/L)$--$\mathcal{O}(1/L^2)$ and its sign is \emph{not} universal ---
it depends on the residual $\SUR$-breaking source (UV boundary conditions, brane
kinetic terms). The clean closed form $-\pi/(4\cwsq L)$ with a definite negative
sign is the \emph{specific minimal-custodial parametrisation our code adopts},
not a theorem; the robust, parametrisation-independent statement is item (1) ---
the net coefficient shrinks by $\sim 1/(2L^2)$. One-loop effects then re-introduce
a sizable, sign-definite contribution --- see \S\ref{sec:loop}.
\item \textbf{$S$ is untouched.} Custodial $\SUR$ does nothing for
$\Pibig'_{33}-\Pibig'_{3Q}$, so \eqref{eq:Smin} stands. This is why the custodial
floor is \emph{$S$-dominated}.
\end{enumerate}
\begin{callout}{Worked example --- one benchmark, all three numbers ($\MKK = 3$~TeV).}
Take $\vev = 246$~GeV, $\MKK = 3000$~GeV, so the decoupling factor is
\[
\frac{\vev^2}{\MKK^2} = \left(\frac{246}{3000}\right)^2 \approx 6.7\times10^{-3}.
\]
Then, with $L=35$, $\cwsq\approx0.769$:
\begin{align*}
\Delta T_{\rm minimal}  &= \tfrac{\pi L}{2\cwsq}\cdot\tfrac{\vev^2}{\MKK^2}
   \approx 71.5\times 6.7\times10^{-3} \approx 0.48
   &&(\text{vs.\ }T=0.047\pm0.066 \Rightarrow \sim 6.6\sigma\text{ in }T\text{ alone},\ \textbf{excluded}),\\
\Delta T_{\rm custodial} &= -\tfrac{\pi}{4\cwsq L}\cdot\tfrac{\vev^2}{\MKK^2}
   \approx -0.029\times 6.7\times10^{-3} \approx -2\times10^{-4}
   &&(\textbf{negligible}),\\
\Delta S &= c_S\cdot\tfrac{\vev^2}{\MKK^2}
   \approx 30\times 6.7\times10^{-3} \approx 0.20
   &&(\text{vs.\ }S=0.026\pm0.075 \Rightarrow \text{already at the edge}).
\end{align*}
This single benchmark contains the entire story: in the minimal model $T$ alone
is a $7\sigma$ disaster (hence $\MKK$ must climb to $\sim 8$--$10$~TeV); custodial
protection makes $T$ vanish, leaving $S\approx0.20$ as the binding constraint ---
which is why the custodial floor sits at a few TeV (at $3$~TeV $S$ is already too
large with $c_S=30$; one needs $\MKK\approx 6$~TeV for the proxy-$S$ alone, see
\S\ref{sec:codescans}).
\end{callout}

\section{$\PLR$: a \emph{separate} protection for $Z\bar b_L b_L$}
\label{sec:PLR}

There is a second, logically distinct, custodial mechanism that is easy to
conflate with the $T$ cure. The non-universal $Z\bar b_L b_L$ vertex correction
(the non-oblique third-generation effect of \S\ref{sec:nonoblique}a) can also be
protected --- but this requires more than $\SUR$. It requires a discrete
left--right parity $\PLR$ (an element exchanging $\SUL \leftrightarrow \SUR$)
under which $b_L$ sits in a representation with $T^3_L = T^3_R$. The
Agashe--Contino--Da Rold--Pomarol argument for why the $Z$ coupling of such a
field is protected is worth internalising, because it is \emph{two constraints
meeting}:
\begin{itemize}[leftmargin=1.4em]
\item gauge invariance under the unbroken $SU(2)_V$ keeps the \emph{sum}
$T^3_L+T^3_R$ fixed, so any vertex shift obeys $\delta T^3_L+\delta T^3_R=0$;
\item $\PLR$ forces left and right to \emph{move together},
$\delta T^3_L=\delta T^3_R$.
\end{itemize}
Two equations force both shifts to vanish, $\delta g_L^b\to 0$. \textbf{In one
line:} \emph{$\PLR$ makes the left and right couplings move together, gauge
invariance forbids the pair from moving, so neither moves.} The protected
quantity is the whole non-universal $\propto(T^3_L-T^3_R)$ vertex correction.

To realise it, the third-generation doublet is embedded as a \textbf{bidoublet}
$Q_L = (t_L, b_L) \in (2,2)_{2/3}$ of $\SUL\times\SUR$ (our code default,
\code{DEFAULT\_CUSTODIAL\_Q\_L\_REP = "all\_gen\_Q\_L\_bidoublet\_(2,2)\_\{2/3\}"}).
Then:
\begin{itemize}[leftmargin=1.4em]
\item the \emph{left-handed} down-type vertex $\delta g_L^b$ is protected by
$\PLR$ (the whole non-universal $Zb_L$ vertex, including the leading gauge piece
$\propto F(c_{Q_3})^2$, not merely the $m_b^2$ admixture);
\item the \emph{right-handed} vertex $\delta g_R^b$ is \emph{not} protected
(there is no $\PLR$ on $b_R$), so $A_b$ / $A_{FB}^{0,b}$ remain a residual,
representation-dependent observable. Our code treats $b_R$ as elementary
(\code{DEFAULT\_CUSTODIAL\_B\_R\_REP = "b\_R\_elementary"},
\code{bR\_strategy = "elementary\_zero"}).
\end{itemize}
\noindent Why single out $A_b$? It is the one $Z$-pole asymmetry that probes the
\emph{right-handed} down-type coupling --- exactly the sector that $\SUR$ and
$\PLR$ leave open. Historically $A_{FB}^{0,b}$ carried a persistent $\sim3\sigma$
tension at LEP/SLD, and an unprotected $b_R$ is precisely the knob that could
shift it. So ``$b_R$ unprotected'' is a \emph{feature}, not just an omission: it
is the one handle on data this class of models retains, which is why our code
(and the literature) keeps it as a separate, representation-dependent residual
instead of zeroing it by symmetry.

\textbf{Two separate clocks.} It is worth stating explicitly, because we got it
wrong once internally (recorded in the orchestration ledger): the $T$-parameter
protection ($\SUR$) and the $Zb_L$ protection ($\PLR$ acting on the bidoublet)
are \emph{different} statements in the \emph{same} custodial group. A model can
have one without the other. In our scan, turning on \code{custodial\_rs\_plr}
flips \emph{both} the $\Delta T$ coefficient \eqref{eq:Tcust} \emph{and} the
$Z\bar b_L b_L$ treatment in \code{T010}/\code{T011}.

\section{Representations and what is/isn't protected}
\label{sec:reps}

The standard custodial embedding (and our code defaults) are:
\begin{center}
\begin{tabular}{llll}
\toprule
Field & Representation & Protection type & Code default \\
\midrule
$Q_L=(t_L,b_L)$ & $(2,2)_{2/3}$ bidoublet & $\PLR$ ($T^3_L=T^3_R$) & \code{...Q\_L\_bidoublet\_(2,2)\_\{2/3\}} \\
$t_R$ & $(1,1)_{2/3}$ singlet & singlet ($T^3_R=0$) & \code{t\_R\_singlet\_(1,1)\_\{2/3\}} \\
$b_R$ & elementary / unprotected & none & \code{b\_R\_elementary} \\
scope & all-gen down-left, diag.\ & --- & \code{all\_gen\_down\_left\_diagonal} \\
\bottomrule
\end{tabular}
\end{center}
Note the table teaches \emph{two distinct} protections: $b_L$ is protected by
the $\PLR$ / vector-like mechanism (which is \emph{why} it needs the bidoublet,
to enforce $T^3_L=T^3_R$), while $t_R$ is protected simply by being an
$SU(2)$ \emph{singlet} ($T^3=0$, so it does not couple to $W^3$ at all). Same
group, different reasons --- which is why $t_R$ sits in $(1,1)$ and $Q_L$ in
$(2,2)$.

Summary of the protection bookkeeping:
\begin{itemize}[leftmargin=1.4em]
\item $\Delta T$: volume-enhanced piece removed by $\SUR$ \eqref{eq:Tcust}. \emph{Protected.}
\item $\delta g_L^b$ ($Z\bar b_L b_L$): removed by $\PLR$ on the bidoublet. \emph{Protected.}
\item $\delta g_R^b$ ($A_b$): \emph{not} protected; residual.
\item $\Delta S$: \emph{not} protected; sets the floor.
\end{itemize}

\begin{callout}{Four common confusions, defused.}
\begin{itemize}[leftmargin=1.2em,nosep]
\item[\textbf{A.}] \emph{``Custodial symmetry protects $T$ \textbf{and}
$Zb_L$.''} No --- continuous $\SUR$ protects $T$; the \emph{discrete} $\PLR$
(requiring the bidoublet, $T^3_L=T^3_R$) protects $Zb_L$. One group, two
different statements; a model can have one without the other.
\item[\textbf{B.}] \emph{``$\MKK$ means one thing.''} It is overloaded by the
factor $x_1\approx2.45$ --- physical first KK mass vs.\ $\LIR$
(\S\ref{sec:Mconvention}). A ``10~TeV'' bound in $\LIR$ units is ``25~TeV''
physical.
\item[\textbf{C.}] \emph{``Custodial also protects $b_R$ / $A_b$.''} No --- $b_R$
is unprotected, which is exactly why $A_{FB}^{0,b}$ stays a live, model-dependent
observable.
\item[\textbf{D.}] \emph{``$S=T=U=0$ means no new physics.''} No --- it is the
\emph{reference SM point} (fixed reference Higgs and top mass); $S,T,U$ measure
deviations \emph{from} that reference, so even the SM itself is not ``at the
origin'' once its inputs differ from the reference.
\end{itemize}
\end{callout}

\section{The calculable one-loop top-partner corrections}
\label{sec:loop}

Here is the deepest payoff of the symmetry, and it is easy to under-sell. The
\emph{same} $\SUR/\PLR$ that kills the tree-level $T$ and $\delta g_L^b$ also
forbids the local 5D counterterms that would otherwise renormalise them; the
symmetry is broken only \emph{non-locally} (by UV-brane boundary conditions, far
from the IR where the loop lives). A would-be UV divergence has no counterterm to
land on, so it cannot be there: the leading \emph{one-loop} contributions from
the light top-partner KK fermions come out \emph{finite and cutoff-independent}.
\textbf{Slogan:} \emph{the symmetry that removes the tree-level effect also
forbids the counterterm that would make the loop ambiguous --- so the loop is a
genuine prediction, not a free parameter.} Moreover the two observables are
\emph{correlated}: they arise from the \emph{same} top-partner multiplets mixing
with the top, so they are not independent numbers but trace a single correlated
band in the $(\Delta T,\ \delta g_L^b)$ plane. A custodial fit constrains them
\emph{jointly}.
Schematically, with $m_{\rm top\text{-}partner}$ the relevant KK fermion mass,
the singlet (custodial partner of $t_R$) gives a \emph{positive}
$\Delta T_{\rm loop}$, while the bidoublet vertex contribution to $\delta g_L^b$
enters with the \emph{opposite} (negative) sign:
\begin{equation}
\Delta T^{\rm singlet}_{\rm loop} > 0, \qquad
\delta g_L^{b,\,\rm bidoublet} < 0,
\label{eq:loopsigns}
\end{equation}
matching the structure of Carena et al., Eqs.~(28)--(30). These one-loop pieces
are $\mathcal{O}(0.1)$ for $T$ and $\mathcal{O}(10^{-3})$ for $\delta g_L^b$ at
$\MKK\sim$~few~TeV --- not negligible at that scale. \textbf{In our code the
hook exists but is deferred by default} (\S\ref{part:code}): the
\code{custodial\_rs\_plr} scan we have run is tree-level, with the loop a flagged,
overridable input (\code{top\_partner\_loop\_delta\_t\_override}).

\section{The bottom line on scales}
\label{sec:scales}

Putting the pieces together (and consistent with the literature and our own
scans):
\begin{center}
\begin{tabular}{lll}
\toprule
Model & Binding EW constraint & Lightest KK gauge mass floor \\
\midrule
Minimal RS (non-custodial) & $\Delta T \propto \pi L/2\cwsq$ & $\sim 8$--$10$~TeV \\
Custodial $\SUL\times\SUR\times\UX\times\PLR$ & $\Delta S$ (T protected) & $\sim 3$~TeV \\
\bottomrule
\end{tabular}
\end{center}
The custodial model trades the volume-enhanced $T$ for an unprotected $S$, which
is numerically much milder, dropping the EW floor from $\sim 8$--$10$ to
$\sim 3$~TeV --- low enough to matter for the LHC, which is the entire
phenomenological motivation. (Compare 0907.0474, which does the analogous thing
for the \emph{flavour} bound by aligning the bulk masses: same spirit, different
sector.)

There is, in fact, a second philosophy for evading the same bound, worth knowing
as a contrast. Instead of adding a \emph{symmetry}, one can change the
\emph{geometry}: Cabrer--von Gersdorff--Quirós deform the AdS metric near the IR
brane (a soft-wall-like deformation) so that the Higgs is softly decoupled from
the KK excitations near the IR --- a strong wavefunction renormalisation that
reduces the Higgs--KK overlap and thereby suppresses \emph{both} $S$ and $T$ ---
with \emph{no} custodial symmetry at all --- reaching $\MKK\sim 1$~TeV. The two
routes,
\[
\textbf{add a symmetry (custodial)} \qquad\text{vs.}\qquad
\textbf{change the geometry (soft wall)},
\]
both confirm the same meta-lesson: the oblique bound is not a fixed number but a
property of the 5D setup. Our scan implements the first route; the second is a
different model (a comparison point, not a drop-in), but it is the cleanest
illustration that ``$S,T$ force $\MKK \gtrsim$ few~TeV'' is a statement about a
\emph{particular} geometry.

% ====================================================================
\part{What is in OUR code}
\label{part:code}
% ====================================================================

\emph{Reader here for the physics may skim Part~\ref{part:code} on a first pass
and jump to the subtleties in Part~\ref{part:subtle}; this part documents the
exact code so the note doubles as an implementation reference.}

\section{The proxy, exactly}
\label{sec:codeproxy}

The oblique sector lives in \code{quarkConstraints/oblique\_stu.py} (physics) with
a thin adapter \code{flavor\_catalog\_constraints/physics\_adapters/oblique\_stu.py}
and the catalogue constraint \code{.../primary/top\_higgs\_ew/EW001.py}. The
predicted shifts are the closed-form proxy
\begin{align}
\Delta S &= c_S\,\frac{\vev^2}{\MKK^2}, & c_S &= 30 \ \ (\text{catalogue anchor}),
\label{eq:codeS}\\
\Delta T_{\rm minimal} &= \frac{\pi L}{2\,\cwsq}\,\frac{\vev^2}{\MKK^2}, &
\Delta T_{\rm custodial} &= -\frac{\pi}{4\,\cwsq\,L}\,\frac{\vev^2}{\MKK^2},
\label{eq:codeT}\\
\Delta U &= 0,
\label{eq:codeU}
\end{align}
with fixed constants (module-level defaults)
\begin{equation}
\vev = 246.21965~\text{GeV},\quad \swsq = 0.23122,\quad L = 35.0,\quad
\chi^2_{95,\,2\rm dof} = 5.991464547107979.
\end{equation}
The $\Delta T$ coefficient functions are \code{minimal\_rs\_t\_coefficient}
($+\pi L/2\cwsq \approx 71.5$) and \code{custodial\_rs\_plr\_t\_coefficient}
($-\pi/(4\cwsq L) \approx -0.029$), dispatched by \code{rs\_oblique\_t\_coefficient}
on the \code{ew\_model} string. The prediction object
\code{rs\_minimal\_oblique\_proxy} assembles
$\Delta T = \Delta T_{\rm tree} + \Delta T_{\rm loop}$, with $\Delta T_{\rm loop}=0$
unless a loop override/metadata is supplied.

\section{The likelihood, exactly}
\label{sec:codelik}

\code{compare\_oblique\_st\_to\_fit} forms the prediction--vs--fit $\chi^2$ of
\eqref{eq:chi2} using the inverse covariance built from
$(\sigma_S,\sigma_T,\rho_{ST})$ in \eqref{eq:fit}, and returns
$(\chi^2,\ \chi^2/5.991,\ \text{budget},\ \text{passes})$. The constraint
\code{EW001} is \code{Severity.HARD}: a point that lands outside the 95\% ellipse
is \emph{vetoed}. The scalar \code{predicted} reported by \code{EW001} is the
$\chi^2$ itself; the vector $(S_p,T_p)$ and the covariance go in
\code{diagnostics}.

\section{The model switch and the deferred loop}
\label{sec:codeswitch}

\code{SUPPORTED\_RS\_EW\_MODELS = ("minimal\_rs", "custodial\_rs\_plr")}; the
default is \code{minimal\_rs}, which is the dataclass default in
\code{oblique\_stu.py}. The stronger ``\emph{byte-identical} to the
pre-custodial code'' guarantee is a property of the \emph{scan-payload} layer,
not of \code{oblique\_stu.py} itself: per the orchestration ledger the
\code{ew\_model} key is popped from the scan config when minimal, so config
hashes are unchanged (ledger hash \code{45e21a07585f7489}). Selecting
\code{custodial\_rs\_plr}:
\begin{itemize}[leftmargin=1.4em]
\item swaps the $\Delta T$ coefficient to \eqref{eq:Tcust};
\item carries the representation metadata of \S\ref{sec:reps}
(\code{qL\_rep}, \code{tR\_rep}, \code{bR\_rep}, \code{protect\_scope},
\code{bR\_strategy});
\item couples to the $Z\bar b_L b_L$ treatment in \code{T010}/\code{T011}
(the $\PLR$ protection there);
\item exposes the one-loop top-partner hook
(\code{include\_top\_partner\_loops}, \code{top\_partner\_loop\_delta\_t\_override})
--- \textbf{off by default}; \code{include\_top\_partner\_loops=True} requires
\code{ew\_model="custodial\_rs\_plr"} (it raises otherwise).
\end{itemize}
The one-loop top-partner quantities exist as named fields in
\code{rs\_ew\_couplings.py} (e.g.\ \code{top\_partner\_delta\_t\_singlet\_magnitude},
\code{top\_partner\_delta\_g\_L\_b\_singlet},
\code{top\_partner\_delta\_g\_L\_b\_bidoublet\_vertex}) but the production scan we
ran set \code{custodial\_top\_partner\_loop\_status = deferred}.

\section{What the scans actually showed}
\label{sec:codescans}

In the paired 1M quark-only scans (baseline job \code{20128400} vs custodial
\code{20675555}, same $r\times\MKK$ grid and seeds):
\begin{itemize}[leftmargin=1.4em]
\item \textbf{Minimal:} the rigorous floor is set by $Z\to b\bar b$
(\code{T010}/\code{T011}), $\sim 25$--$30$~TeV in \emph{physical} $\MKK$
($\approx 10$--$12$~TeV in $\LIR$ units; \S\ref{sec:Mconvention}). The proxy
oblique \code{EW001} and $\Delta F=2$ are subdominant to $Zb\bar b$ here.
\item \textbf{Custodial:} \code{T010} is removed by $\PLR$ (verified:
$694{,}123$ minimal vetoes $\to 0$ custodial). The \emph{rigorous} floor drops to
the $\Delta F=2$ set ($\eps_K$/$\Delta m_s$/$\phi_s$ = \code{K001}/\code{B003}/\code{B004}),
$\sim 2$--$3$~TeV. The \emph{inclusive} floor ($\sim 7$~TeV) is then set by the
\emph{proxy} \code{EW001} oblique-$S$ plus the collider \code{CR*} searches.
\end{itemize}
\textbf{Important honesty point:} because \code{EW001} is tagged \emph{proxy}, it
does \emph{not} count toward the ``rigorous'' floor. So the headline ``custodial
$\to 2$--$3$~TeV'' is the rigorous-only number; the proxy $S$ (with $c_S=30$)
actually still bites at $\sim 6$~TeV. The gap between $2$--$3$ and $\sim 7$~TeV is
\emph{entirely} the not-yet-rigorous oblique-$S$ and collider proxies --- which is
the strongest practical reason to make $S$ rigorous.

% ====================================================================
\part{Subtleties and the gap to ``rigorous''}
\label{part:subtle}
% ====================================================================

\section{The proxy is parametric, not spectral}
The coefficients $c_S=30$ and $c_T = \pm$(\,$\pi L/2\cwsq$ or $\pi/4\cwsq L$\,)
are \emph{closed-form parametric estimates}, not computed per point from the KK
spectrum and wavefunction overlaps. The design-consensus note planned to replace
them with the explicit leading KK-gauge-sum from \code{rs\_ew\_spectrum} (the same
machinery that already feeds $Zb\bar b$), but \code{EW001} still uses the
constants. This is why it is tagged \code{PARTIAL}/\code{NEEDS-HUMAN-PHYSICS},
not \code{rigorous}.

\section{The $c_S=30$ number}
$c_S=30$ comes from the \code{EW001.yaml} ``warped-extra-dimension S
coefficient'' anchor; the YAML also records a one-sided $S\lesssim 0.18$ and a
$\MKK \gtrsim 3.2$~TeV context value. An $\mathcal{O}(30)$ coefficient is on the
large/conservative side --- canonical tree-level RS estimates give $c_S$ closer
to $2\pi$--$\mathcal{O}(10)$ depending on the gauge content and on UV-brane
kinetic terms. \textbf{This single number drives the custodial inclusive floor},
so its provenance and convention (and the associated $\MKK$ definition) are worth
nailing down before quoting a custodial $S$-bound.

\section{$M_{\rm KK}$ convention: the factor of $2.45$}
\label{sec:Mconvention}
``$\MKK$'' is overloaded. The physical first KK gauge mass is
$\MKK^{\rm phys} = x_1\LIR \approx 2.45\,\LIR$ \eqref{eq:KKmass}. Literature
oblique formulas are sometimes quoted with $\MKK \equiv \LIR$. A bound stated as
``$\sim 10$~TeV'' in $\LIR$ units is ``$\sim 25$~TeV'' in physical units --- a
$2.45\times$ difference. Our $Zb\bar b$ note already flagged this for the
$Z\to b\bar b$ floor; the \emph{same} care is needed for the oblique
$\vev^2/\MKK^2$ in \eqref{eq:codeS}--\eqref{eq:codeT}. Whichever convention
$c_S=30$ was calibrated in must match the $\MKK$ fed to \code{EW001}.

\section{Brane kinetic terms (BKTs)}
UV-brane-localised gauge kinetic terms are free parameters that shift both $S$
and $T$ and can be used to \emph{relax} the bounds further (this is part of how
some non-custodial constructions survive). Our proxy sets them to zero; they are
flagged as an explicit human model choice and are not scanned.

\section{Non-oblique $W,Y$ and the projection onto two numbers}
As in \S\ref{sec:nonoblique}, RS is not purely oblique. We keep only $(S,T)$ with
$U=0$ and drop $W,Y$ and the generation-dependent vertex corrections (the latter
are partly captured elsewhere, e.g.\ $Zb\bar b$ in \code{T010}/\code{T011}, but
not folded into a single EW likelihood). A fully rigorous treatment would either
use $\{\hat S,\hat T,W,Y\}$ or fit directly to $Z$-pole + $M_W$ observables.

\section{The one-loop custodial corrections are off}
The calculable, correlated one-loop top-partner contributions to $T$ and
$\delta g_L^b$ (\S\ref{sec:loop}) are coded but deferred in the scan. At
$\MKK\sim$~few~TeV they are not negligible ($\Delta T_{\rm loop}\sim 0.1$), and a
loop-on custodial rerun would shift the custodial floors somewhat.

\section{Checklist: what ``rigorous $S,T,U$'' would require}
\begin{enumerate}[leftmargin=1.8em]
\item Replace \eqref{eq:codeS}--\eqref{eq:codeT} with a per-point KK-gauge-sum
$S,T$ from \code{rs\_ew\_spectrum} (vector profiles + tower), in a fixed,
documented $\MKK$ convention.
\item Decide and implement a BKT policy (zero, or scanned with a prior).
\item Turn on the one-loop top-partner $T$/$\delta g_L^b$ (PR2 numerics) for the
custodial mode.
\item Either add $W,Y$ or move to a direct-observable EW fit, since the
two-parameter projection is itself an approximation in the light-KK regime.
\item Reconcile the $c_S$ normalisation against a cited spectral computation
(e.g.\ Casagrande et al.\ 1005.4315).
\end{enumerate}

% ====================================================================
\part*{Annotated reading guide}
\addcontentsline{toc}{section}{Annotated reading guide}
% ====================================================================

\begin{description}[leftmargin=0em,style=nextline]
\item[Peskin \& Takeuchi, \emph{PRD 46 (1992) 381}.] The original $S,T,U$
definitions \eqref{eq:Tdef}--\eqref{eq:Udef}. Read \S\ref{sec:STUdef} alongside
it. Foundational, not RS-specific.

\item[Barbieri, Pomarol, Rattazzi, Strumia, hep-ph/0405040.] Why $S,T,U$ is
incomplete for light new states; defines the extended set $\hat S,\hat T,W,Y$.
The reference for \S\ref{sec:nonoblique}/\S\ref{part:subtle} ($W,Y$ caveat).

\item[Agashe, Delgado, May, Sundrum, hep-ph/0308036.] \emph{The} custodial-$T$
paper: introduces $\SUL\times\SUR$ in the RS bulk to kill the volume-enhanced
$T$. The holographic ``composite Higgs lacks custodial isospin'' picture of
\S\ref{sec:rsorigin}/\S\ref{sec:Tcure} is theirs.

\item[Agashe, Contino, Da Rold, Pomarol, hep-ph/0605341.] The $\PLR$ protection
of $Z\bar b_L b_L$ (\S\ref{sec:PLR}); the bidoublet embedding and the
$T^3_L=T^3_R$ argument. Cited directly in our \code{T010}/\code{T011}.

\item[Carena, Pontón, Santiago, Wagner, hep-ph/0607106 \& hep-ph/0701055.]
Light-KK custodial RS with explicit oblique fits; \textbf{0701055} has the
one-loop top-partner $T$ and $\delta g_L^b$ (Eqs.~28--30) that our PR2 implements
(\S\ref{sec:loop}). The ``$\sim 3$~TeV'' custodial floor traces to these.

\item[Casagrande, Goertz, Haisch, Neubert, Pfoh, 1005.4315
(\emph{JHEP 09(2010)014}).] The most explicit, pedagogical worked treatment:
mass-basis KK decomposition without tower truncation, closed-form $T$ and
$Zb_L$-protection formulas, minimal ($\sim 8$~TeV) vs custodial ($\sim 3$~TeV)
bounds. \textbf{This is the reference to mine if/when we upgrade \code{EW001}
from proxy to spectral.} Their flavour companion 0807.4937 is the one our
$Zb\bar b$/$\Delta F=2$ code already cross-checks against.

\item[Cabrer, von Gersdorff, Quirós, 1103.1388 (\& 1011.2205, JHEP 10(2011)137).]
The \emph{alternative} cure: deform the IR metric (soft-wall-like) to suppress
$S$ and $T$ \emph{without} custodial symmetry, reaching $\MKK\sim 1$~TeV. A
different model than ours --- a comparison point, not a drop-in --- but the
clearest illustration that the oblique bound is geometry-dependent. Spiritually
closest to 0907.0474 (change the setup to evade the bound).

\item[``Bulk RS, EW precision tests and the 125 GeV Higgs'', \emph{PRD 93
(2016) 075008}.] Post-Higgs redo of the $S,T$ fit with the measured Higgs mass;
useful for current numbers.
\end{description}

\bigskip
\noindent\rule{\linewidth}{0.4pt}\\
\footnotesize
\emph{Internal note. Equations \eqref{eq:Tmin}, \eqref{eq:Smin}, \eqref{eq:Tcust}
are the parametric forms our code uses; the numerical coefficients
($c_S=30$, $L=35$, $c_T=\pm$) and conventions ($\MKK$, $\vev=246$~GeV) are as
stored in \code{oblique\_stu.py} and \code{EW001.yaml}. All ``rigorous vs proxy''
tags follow the orchestration ledgers. Cross-check before citing externally.}

\end{document}
